% Scope: Synthesize sequence and parameter choices by clinical objective, with anatomy used for examples rather than the book's organizing principle.
% Status: architecture only; no chapter prose or completed problems yet.
% Prerequisites: Chapters 7 through 11 and 19 through 22.
% Planned worked examples: Brain FLAIR, cardiac cine/LGE and liver fat-water protocol comparisons.
% Planned figures: Protocol decision maps tied to physical contrast.
% Planned challenge problems: Design a fixed-time protocol; justify every trade-off and identify failure modes.
\chapter{Clinical Protocol Physics}
\label{ch:23}

\section{From diagnostic task to acquisition}
\label{sec:23:from-diagnostic-task-to-acquisition}
\subsection{Target contrast and spatial scale}
\subsection{Motion, timing and patient constraints}
\subsection{Redundancy, uncertainty and protocol robustness}

\section{Brain protocols}
\label{sec:23:brain-protocols}
\subsection{Structural MPRAGE and FLAIR suppression}
\subsection{Acute ischemia and diffusion interpretation}
\subsection{Susceptibility, hemorrhage and T2-star}

\section{Cardiac protocols}
\label{sec:23:cardiac-protocols}
\subsection{bSSFP cine and off-resonance management}
\subsection{Gating, motion and temporal resolution}
\subsection{Late gadolinium enhancement and inversion timing}

\section{Liver and abdominal protocols}
\label{sec:23:liver-and-abdominal-protocols}
\subsection{Breath holding and respiratory motion}
\subsection{Dixon, fat fraction and iron-sensitive decay}
\subsection{Diffusion and dynamic contrast timing}

\section{Musculoskeletal and oncologic protocols}
\label{sec:23:musculoskeletal-and-oncologic-protocols}
\subsection{Fat suppression and fluid-sensitive imaging}
\subsection{Spatial resolution and partial-volume control}
\subsection{Multiparametric tumor assessment and specificity}

\section{Vascular protocol synthesis}
\label{sec:23:vascular-protocol-synthesis}
\subsection{TOF, phase contrast and contrast enhancement}
\subsection{Velocity range and flow artifacts}
\subsection{Choosing coverage, resolution and scan time}

\section{Worked examples}
\label{sec:23:worked-examples}
% Planned calculations: Brain FLAIR, cardiac cine/LGE and liver fat-water protocol comparisons.
\subsection{Derivation and quantitative calculation}
\subsection{Assumptions, limiting cases and scanner interpretation}

\section{Challenge problems}
\label{sec:23:challenge-problems}
% Problem design: Design a fixed-time protocol; justify every trade-off and identify failure modes.
% Use challengeproblem and stable labels prob:23:<topic>.
% Complete solutions belong in Appendix E, section for this chapter.
% Use \begin{solution}{prob:23:<topic>} ... \end{solution} there.
\subsection{Derivation and quantitative design}
\subsection{Model criticism and integrated reasoning}
